% Fichier généré automatiquement par tools/generate_corrections.py.
% Source : DYN/DYN-03-Inertie/1026_Pale/1026_Pale.tex
% Statut : complet (6/6 question(s) renseignée(s), 0 reprise(s)).
\begin{corrigebox}[Corrigé extrait de la banque]
\CorrigeQuestion{1}
La portion 1 est un cylindre $\vect{OG_1} = -\dfrac{\ell}{2}\vx{}$.
\CorrigeQuestion{2}
$m_1 =\mu \pi r^2 \ell$

$m_2 = \mu \dfrac{1}{2} \pi R^2 L + \mu RHL$
\CorrigeQuestion{3}
La portion $2_a$ est composée d'un demi cylindre de centre d'inertie $G_a$. Pour des raisons de symétrie, seule la composante suivant $\vz{}$ est à déterminer.

On a donc $m_a \vect{OG_a}\cdot \vz{} = \iint \vect{OP}\cdot \vz{} \d m$ 
$ =\iint \rho \sin \theta \mu      \rho \d\rho \d \theta \d z$  
$ = - \mu    \dfrac{1}{3}R^3  L \left[ \cos \theta \right]_{0}^{\pi}  $  
$ =  2 \mu    \dfrac{1}{3}R^3  L $

Avec $m_a = \mu \pi R^2/2L$, on a $ \vect{OG_a}\cdot \vz{} = \dfrac{4}{3\pi}R  $.

On a donc $\vect{OG_a}=-\left(\ell+\dfrac{L}{2}\right) \vx{}+\dfrac{4}{3\pi}R \vz{}$.

La portion $2_b$ es un prisme. On a $\vect{OG_b}=-\left(\ell+\dfrac{L}{2}\right) \vx{}-\dfrac{1}{3}H \vz{}$. Elle est de masse $m_b = \mu RHL$.

On a alors $\vect{OG_2} = \dfrac{m_a \vect{OG_a} + m_b \vect{OG_b} }{m_2}$
\CorrigeQuestion{4}
$\vect{OG} = \dfrac{m_1 \vect{OG_1} + m_2 \vect{OG_2} }{m_1+m_2} =-a\vx{}-c\vz{}$ .
\CorrigeQuestion{5}
En $O$ il y a une inifinité de plans de symétrie pour 1. De plus, $\vy{}$ et $\vz{}$ jouent le même rôle. 
$\inertie{O}{1} = \matinertie{A_1}{B_1}{B_1}{0}{0}{0}{\rep{}}$.

En $G_2$ il y a 2 plans de symétrie perpendiculaires pour 2 $(G_2\,\vy{},\vz{})$ et $(G_2\,\vz{},\vx{})$ .
$\inertie{G_2}{2} = \matinertie{A_2}{B_2}{C_2}{0}{0}{0}{\rep{}}$.

En $O$ il y a 1 plans de symétrie permendiculaires pour 1+2  $(G_2\,\vz{},\vx{})$ .
$\inertie{O}{1+2} = \matinertie{A_{1+2}}{B_{1+2}}{C_{1+2}}{0}{-E_{1+2}}{0}{\rep{}}$.
\CorrigeQuestion{6}
On a  $\vect{OG_2} = a\vx{}+c\vz{}$ et 
$\inertie{O}{2} = \inertie{G_2}{2} + \matinertie{m_2c^2}{m_2 \left(a^2+c^2\right)}{m_2a^2}{0}{-m_2 ac }{0}{}$

$=\matinertie{A_2+m_2c^2}{B_2 + m_2 \left(a^2+c^2\right)}{C_2 + m_2a^2}{0}{-m_2 ac }{0}{}$

$\inertie{O}{1+2} =\matinertie{A_1 + A_2+m_2c^2}{B_1 + B_2 + m_2 \left(a^2+c^2\right)}{C_1 + C_2 + m_2a^2}{0}{-m_2 ac }{0}{}$
\end{corrigebox}
